\documentclass{article}
\usepackage{amsmath,amssymb}
\usepackage{xurl}
\usepackage[hidelinks]{hyperref}
% Shared commands for notes in docs/paper-gaps/.

\DeclareUnicodeCharacter{2080}{\ifmmode{}_{0}\else\textsubscript{0}\fi}
\DeclareUnicodeCharacter{2081}{\ifmmode{}_{1}\else\textsubscript{1}\fi}
\DeclareUnicodeCharacter{2082}{\ifmmode{}_{2}\else\textsubscript{2}\fi}
\DeclareUnicodeCharacter{2099}{\ifmmode{}_{n}\else\textsubscript{n}\fi}

\newcommand{\Lean}{\textsc{Lean}}
\ExplSyntaxOn
\NewDocumentCommand{\leanid}{m}{
  \ifmmode
    \textsf{\detokenize{#1}}
  \else
    \group_begin:
    \urlstyle{sf}
    \tl_set:Nn \l_tmpa_tl {#1}
    \tl_replace_all:Nnn \l_tmpa_tl {\_} {_}
    \exp_args:NV \path \l_tmpa_tl
    \group_end:
  \fi
}
\ExplSyntaxOff
\providecommand{\tr}{\operatorname{tr}}
\newcommand{\tnleanIssueBase}{https://github.com/LionSR/TNLean/issues/}
\newcommand{\tnleanPrBase}{https://github.com/LionSR/TNLean/pull/}
\newcommand{\tnleanDocBase}{https://sirui-lu.com/TNLean/docs/}
\newcommand{\ghissue}[1]{\href{\tnleanIssueBase#1}{GitHub issue \##1}}
\newcommand{\ghpr}[1]{\href{\tnleanPrBase#1}{PR \##1}}
\newcommand{\doclink}[1]{\href{\tnleanDocBase#1.html}{\path{#1}}}

% Dirac notation and the identity operator, as in blueprint/src/macros/common.tex.
% \providecommand keeps note-local definitions authoritative.
\usepackage{dsfont}
\providecommand{\ket}[1]{|#1\rangle}
\providecommand{\bra}[1]{\langle #1|}
\providecommand{\braket}[2]{\langle #1 | #2 \rangle}
\providecommand{\ketbra}[2]{|#1\rangle\!\langle #2|}
\providecommand{\Id}{\mathds{1}}
\providecommand{\one}{\mathds{1}}
\providecommand{\C}{\mathbb{C}}
\providecommand{\MN}[1]{M_{#1}(\C)}
\providecommand{\MD}{M_D(\C)}

% Verdict marker: \gapnote{<kind>}{<status>}. Kind is the policy
% classification (clarification, local-correction, scope-restriction,
% unfaithful, false-source, open-gap); status is open, wip (elimination
% underway), resolved, or historical. Machine-read by the site index;
% prints a verdict line.
\newcommand{\gapnote}[2]{\par\noindent\textbf{Verdict:} #1 (#2).\par}

\title{Compatible Coherent Seeds and Non-Normal State Conversion}
\author{}
\date{2026-10-05}
\begin{document}
\maketitle
\gapnote{scope-restriction}{open}

\section{The unresolved interpretation}
The discussion of~\cite{Malz2024Preparation} states that MPS in the same
phase can be transformed into each other with a logarithmic-depth circuit.
It does not specify a non-normal phase predicate or how a selected vector
in a degenerate groundspace is transported. Issue
\href{https://github.com/LionSR/TNLean/issues/8469}{8469} asks for a circuit
between normalized periodic vectors, without a compatibility condition on
their sector amplitudes. The obstruction below concerns the interpretation
that equality of parent groundspaces suffices for this chosen-vector claim.
It does not refute every interpretation of the paper's informal remark.

\section{A quantitative obstruction}
Consider the tensors
\begin{align}
    A^0 & =\operatorname{diag}(1,1,0), & A^1 & =\operatorname{diag}(0,0,1), \\
    B^0 & =\operatorname{diag}(1,0),   & B^1 & =\operatorname{diag}(0,1).
\end{align}
For every positive length their normalized periodic vectors are
\begin{align}
    \psi_N & =(2|0^N\rangle+|1^N\rangle)/\sqrt5, &
    \phi_N & =(|0^N\rangle+|1^N\rangle)/\sqrt2.
\end{align}
Both have local parent support
$\operatorname{span}\{|00\rangle,|11\rangle\}$ and the same ferromagnetic
parent groundspace. All copy weights are nonzero. The first tensor and its
periodic-vector formula already occur in
\doclink{TNLean/MPS/Preparation/RepeatedBlockCounterexample}; no claim about
its previously corrected polar approximation is needed here.

Let $U$ be a depth-$T$ nearest-neighbor unitary circuit on the ring, and
assume $N>4T+4$. Choose two sites whose backward light cones are disjoint
and whose union leaves at least one site unobserved. The pulled-back Pauli
observables $X=U^\dagger Z_iU$ and $Y=U^\dagger Z_jU$ have norm one and
are supported in those respective cones. Their cross terms between
$|0^N\rangle$ and $|1^N\rangle$ vanish, as do those of $XY$, because of
the unobserved site. Expectations of $XY$ factor within each product
branch. Writing $p=4/5$, $a_k=\langle k^N|X|k^N\rangle$ and
$b_k=\langle k^N|Y|k^N\rangle$, with $a_k,b_k\in[-1,1]$, gives
\begin{align}
    \operatorname{Cov}_{U\psi_N}(Z_i,Z_j)
     & =p(1-p)(a_0-a_1)(b_0-b_1)\leq 16/25.
\end{align}
The covariance in $\phi_N$ equals one. With the convention
$D(\rho,\sigma)=\tfrac12\|\rho-\sigma\|_1$, every norm-one observable
has expectation difference at most $2D$. The covariance difference is
therefore at most $6D$, so $D\geq3/50$. For pure states, putting
$F=|\langle\phi_N|U\psi_N\rangle|$, one obtains
\begin{align}
    D^2=1-F^2\leq2(1-F),\qquad 1-F\geq9/5000.
\end{align}
Fixing an accuracy below this bound contradicts any depth
$T\leq c\log(N/\varepsilon)$ for all sufficiently large $N$.
Local product ancillas do not remove the product-branch factorization;
measurements with nonlocal classical feedforward are a different resource.

\paragraph{Proof status.}
The covariance obstruction in this section is a mathematical argument,
not a Lean-checked theorem. The remaining formalization is tracked by
issue~8469. In particular, the existing multiplicity counterexample does
not itself prove this circuit lower bound.

\section{The sufficient statement}
Suppose normalized vectors $a,b$ are approximated by $U_A\chi_A,U_B\chi_B$
with errors $\varepsilon_A,\varepsilon_B$, where both seeds are normalized
and $R\chi_A=\chi_B$. Require $U_A,U_B,R$ to be actual circuits on one
common physical ring, of depths $T_A,T_B,T_R$. Then
\begin{align}
    U                       & =U_BRU_A^\dagger,                   & T_U & \leq T_A+T_R+T_B, \\
    1-|\langle Ua|b\rangle| & \leq2(\varepsilon_A+\varepsilon_B).
\end{align}
The common-seed case has $R=I$, $T_R=0$. Taking endpoint errors at most
$\varepsilon/4$ and uniform depth bounds yields $O(\log(N/\varepsilon))$.
The constants must precede both $N$ and $\varepsilon$ in the quantifiers.
No symmetry constraint is imposed, and no constant-depth or exact target
conversion is claimed. Any auxiliary registers must be included with their
specified reset vectors; reduced-state equality is insufficient to invert
a preparation.

The coherent circuit extracted from the non-normal preparation implements
all encoded sector basis vectors simultaneously; its choice does not depend
on their amplitudes. Its input includes zero-valued unused sites and its
output is the full block-isometry vector. Its depth is at most $CL$ for
physical block lengths bounded by $L$. An arbitrary growing-block unitary
must not be counted as a free matching gate.

For GHZ seeds a permutation of sector labels and matching amplitude
magnitudes allow relative phases to be corrected on a single encoded
register. This observation is not used as an unproved circuit assumption:
the conversion theorem requires the actual circuit $R$ and counts its
physical-site depth. General canonical data do not imply that matching.
Complex weights, multiplicities and length dependence must be retained;
when using the corrected one-copy construction, its corrected input
amplitudes must be used, not silently replaced by the printed coefficients.

\bibliographystyle{alpha}
\bibliography{references}
\end{document}
